\zlabel{triptych:concise:commentary:start}
\section{Preparing for a meeting that cannot be postponed}
\subsection{What the lamps do not establish}
The parable begins with people already on their way to meet the bridegroom. Their division does not separate those who expect him from those who have never heard of him. The ten share virginity, lamps, the delay, sleep, and the awakening cry. Five have also brought oil in vessels. Before explaining the symbols, the narrative requires this distinction between a common outward preparation and a provision that some have neglected.

Chrysostom takes the choice of virgins seriously as a choice of people who have undertaken difficult self-denial. Their failure cannot be explained by saying that they have made no effort. He compares the struggle against bodily appetite with attachment to money: a person has overcome the former yet yields to the latter. The oil's identification with mercy therefore does more than add generosity to a list of virtues. It exposes a neglected obligation that an admired achievement has allowed the person to overlook. The neighbor remains hungry regardless of how demanding the giver's other sacrifices have been.\footnote{Chrysostom, \textit{Homilies on Matthew}, 78.1.}

Augustine starts from the same common equipment but gives the lamps a different scope. They are works that others can see. His examples include feeding, clothing, visiting, welcoming, reconciling and burying the dead. Some of these are precisely the acts whose absence troubles Chrysostom. Augustine's foolish virgins can perform them. He explains virginity more broadly too, as restraint of the senses from unlawful indulgence, so that the admonition reaches married Christians living faithfully as well as consecrated virgins. The Gospel questions the purpose of recognizably good conduct within the Church.\footnote{Augustine, Sermon 93, §§1--5.}

Augustine explains the possibility of such conduct without charity through Paul's account of love. He compares oil rising above other liquids with charity's preeminence. An act can benefit the receiver and yet be performed for the giver's reputation. That distinction does not erase the benefit: the hungry person has actually been fed. It explains why an observer's accurate report of the deed cannot establish the purpose from which it came. The person can need conversion even while continuing an action that ought to be done.

These interpretations consequently direct examination toward different facts. Chrysostom asks about help refused despite the means to give it. Augustine asks about the love within help already given. Treating both explanations as an instruction simply to do more would miss Augustine's question; treating both as an inquiry into feelings would miss Chrysostom's neglected person. Discovery of vanity is a reason to seek charity in continuing the service. It supplies no reason to suspend assistance until the giver feels free of every mixed motive.

\subsection{Why the wise cannot provide what is missing}
The request for oil tests Chrysostom's interpretation at its most vulnerable point. The foolish ask for a share; the need is urgent, and the wise have a supply. A refusal here appears incompatible with identifying oil as mercy. Chrysostom presents the difficulty himself. His explanation depends on the time represented by the arrival: one person's works cannot become another's after death. The wise describe an inability to transfer a life already lived. The parable's final separation differs from a present encounter in which food, money or attention can still be given.\footnote{Chrysostom, \textit{Homilies on Matthew}, 78.1--2.}

This temporal distinction also explains his sellers. They are the poor, available during the life in which mercy could have been practiced. The foolish should have approached them earlier. Chrysostom presses the point against delaying generosity until a deathbed disposition of property. Such a plan assumes that time and capacity for that last act will remain. Sudden death can remove both. His instruction is to use the present opportunity, rather than to treat the prospect of a later bequest as if the help had already been given.

The buying image concerns the lasting consequence of generosity. The poor do not sell grace, and gifts create no debt that obliges God to admit the giver. Likewise, the parable's refusal gives no permission to abandon the dying. They are living persons who may still receive assistance and comfort. Chrysostom addresses the hearer's postponement; his explanation supplies no verdict about the fate of a particular deceased person.

Augustine and Gregory explain the request through testimony. In Gregory's account, the foolish ask others to attest the good actions they have seen. The vessels signify hearts, and human favor has been the reward sought by works that looked admirable. A companion can describe a deed; a companion cannot supply the inward charity with which another person should have done it. God already knows the heart. External testimony cannot conceal a purpose from him.\footnote{Gregory, \textit{Homiliae in Evangelia}, 12.1--3; Augustine, Sermon 93, §§9--14.}

Their identification of the sellers as flatterers follows from that account. Praise sustained the foolish before the cry, and they seek its support again. Augustine treats the direction to buy as a rebuke exposing the uselessness of the old dependence. In Chrysostom, the sellers signify persons who should have received help; in Augustine and Gregory, they signify persons from whom approval had been sought. Combining these into a single identification would change the explanation of the failure.

The wise virgins' concern that there might not be enough also requires care. Augustine finds humility in it: a good conscience still seeks forgiveness before God's deeper judgment. Gregory likewise says that one's own testimony scarcely suffices for oneself. The distinction between inward charity and public reputation therefore does not make private satisfaction an assurance of salvation. The plea for merciful hearing at the entrance remains fitting for someone who has tried to act faithfully.

\subsection{A gift that becomes the seeker's own}
The requirement to bring oil can sound like a demand to produce goodness independently. Augustine addresses that difficulty with the olive: the person can put oil in a vessel but did not create its source. Inward charity is received from God. Responsibility concerns the gift's real presence and use within the person. It does not transfer the origin of grace to human achievement.\footnote{Augustine, Sermon 93, §§9, 14.}

Wisdom 6:12--16, corresponding to Vulgate 6:13--17, describes this relation through several actions. The person loves, seeks, rises early and thinks; Wisdom shows herself, anticipates desire, waits nearby and goes out to meet. The historical study English's ``preventeth'' means that she comes before the seeker. Bonaventure explains her visibility through her own light. Intellect and will open to her, while love and action belong to seeking. The search does not begin with a capacity wholly independent of what it seeks.\footnote{Bonaventure, \textit{Commentarius in librum Sapientiae}, chapter 6, Vulgate vv. 13--17, Quaracchi VI, pp. 147--149.}

Wisdom at the door and the closed wedding door describe different circumstances. Wisdom's approach encourages a willing response; Jesus' parable warns that preparation has an end beyond which it cannot be improvised. Read together, they give a reason for confidence now and urgency about using the opportunity. Past failure need not make a new beginning seem inaccessible. Nor does the nearness of the gift make responding unnecessary.

The Collect's petition for freedom to serve can be applied to both difficulties. Attachment to possessions can keep help indefinitely postponed. Dependence on praise can restrict help to circumstances in which it will be noticed. These are different impediments, and relief from one does not establish relief from the other. In either case, the prayer concerns the service of the whole person. Bodily and mental burdens remain real objects of petition; the inward application does not explain every hardship as a failure of intention.

Psalm 62 Vulgate, 63 Hebrew, describes desire and remembrance in the body. Augustine interprets its bodily thirst in relation to resurrection and distinguishes a wakeful soul from a sleepless body. The hands raised in prayer also extend in mercy. Thus an appeal to inward charity remains answerable to action toward a neighbor, while bodily weariness gives no proof that charity has failed.\footnote{Augustine, \textit{Expositions on the Psalms}, 62.1--12, especially §§3--5, 10.}

The entrance lament supplies a related qualification. Augustine reads Psalm 87 in Christ's voluntarily endured human suffering. Divine hearing means compassionate acceptance; sorrow need not mean loss of grace. A person who seeks a more truthful purpose need not manufacture a cheerful appearance as evidence of success. Bellarmine's exposition of Psalm 62 likewise explains seeking in dependence on help already received. Desire can continue alongside recollection of that help.\footnote{Augustine, \textit{Expositions on the Psalms}, 87.1--4; Bellarmine, \textit{Commentary on the Psalms}, O'Sullivan abridgment, Psalm 62:1--8.}

\subsection{Sleep, the cry and the grief of those who wait}
Augustine tests a possible meaning of sleep against the complete parable. If it signifies charity cooling, the wise lose charity as well, because all sleep. Their entrance would then contradict that explanation. He instead identifies the common sleep with death and relates it to Paul's language concerning the dead. Watchfulness belongs to faith, hope, charity and works during life. It is compatible with the ordinary rest of the body and does not make exhaustion the reason for exclusion.\footnote{Augustine, Sermon 93, §§6--9, 17.}

Chrysostom also reads the sleep as death. He takes the bridegroom's delay against the expectation of an immediate arrival and connects the midnight cry with Paul's description of the Lord's descent. On midnight, however, the explanations diverge. Chrysostom allows a literal nocturnal inference about the resurrection; Augustine emphasizes unexpected arrival and rejects successive calculations of the end. Neither provides a usable future date. The acclamation's ignorance of the hour calls for continuing service with what has already been entrusted.\footnote{Chrysostom, \textit{Homilies on Matthew}, 78.1 and 77.1--3; Augustine, Sermon 93, §8; Matthew 24:42a, 44.}

Paul's concern in 1 Thessalonians 4:13--18 is the place of the dead in Christ's coming. The living gain no advantage that excludes them. Chrysostom explains the meeting through the welcome of an arriving king and distinguishes resurrection from the glory of the final meeting. The letter's consolation rests on Christ's death and rising before it describes the gathering. The long form ends with a direction to comfort one another, making this teaching something believers offer to those who mourn.\footnote{Chrysostom, \textit{Homilies on First Thessalonians}, 6--8, opening pericope expositions.}

Aquinas specifies why hope leaves room for sorrow: death causes bodily loss and separation. He explains Christ's resurrection as both cause and pattern of the faithful person's resurrection. The promise is therefore independent of the mourner's success in appearing composed. Chrysostom's exhortations against lament are markedly severe; Aquinas's explicit allowance should not be read into every phrase of them. Their common resurrection hope does not make their pastoral language about grief identical.\footnote{Aquinas, \textit{Super I ad Thessalonicenses}, 4, lectio 2.}

The short form, 4:13--14, retains the foundation in Jesus' death and resurrection. The trumpet, rising order, gathering and final direction to console belong to its continuation. Both forms can sustain the mercy owed to someone grieving; only the long form proclaims the full scene. In the first interpretation, consolation is an act of care that attends to the loss. In the second, sorrow is compatible with inward trust even when it offers no reassuring appearance to spectators. Both anticipate a bodily future with Christ, rather than the mere endurance of a favorable memory.

\subsection{Recognizing the guest and receiving his care}
The Luke Communion alternative refers to the disciples' report in Jerusalem of recognition at Emmaus. Gregory follows the earlier encounter: the disciples discuss Jesus with love and doubt, hear the Scriptures explained, and urge the traveler to stay. Recognition follows their welcome at the meal. They receive someone whose identity they have not yet grasped. Gregory asks whether those unable yet to love him as God will love him as a stranger. Doing what the word teaches becomes a way into understanding it.\footnote{Gregory, \textit{Homiliae in Evangelia}, 23.1--2, on Luke 24:13--35.}

Gregory's account contributes something beyond a second endorsement of generosity. Chrysostom's virgins must have practiced mercy before the final meeting. Gregory explains a present encounter in which hospitality precedes recognition. The disciples do not wait for certainty about the guest before offering help. Together, these arguments support an interpretation in which preparation to meet Christ includes receiving the person whose need can already be met.

The same hospitality tests the inward interpretation. The disciples' love and doubt are expressed through what they do. Their invitation does not wait for an imagined perfection of interior understanding. In Gregory's account, action illuminates understanding. Augustine's concern with purpose can therefore accompany practical welcome without making the neighbor wait while the giver inspects the heart. Aquinas similarly distinguishes the inward defect in the virgins from negligence of action in the following parable of the talents. That following parable supplies context, not an additional appointed Gospel.\footnote{Aquinas, \textit{Super Matthaeum}, chapter 25, lectio 1 and opening of lectio 2.}

The psalm Communion alternative begins with receiving rather than hosting. Only the adaptation of Psalm 22:1--2 Vulgate, 23:1--2 Hebrew, is appointed. Augustine explains Christ feeding the Church and leading to baptismal refreshment; the rest of the psalm gives the wider setting of danger, table and abiding home. Bellarmine's explanation allows hardship to persist under the shepherd's care. Consequently, the inference from receiving care to offering it does not require prosperity as its proof. Someone still in need can both receive assistance and assist another.\footnote{Augustine, \textit{Expositions on the Psalms}, 22.1--6; Bellarmine, \textit{Commentary on the Psalms}, O'Sullivan abridgment, Psalm 22.}

These are alternatives at Communion. The shepherd's feeding and the disciples' welcome remain distinct acts, with different directions of giving and receiving. Either can illuminate dependence on Christ. The Luke passage adds the particular sequence by which hospitality accompanies recognition, while the psalm begins with the Lord's provision.

\subsection{Being known and continuing in charity}
The bridegroom's denial of knowledge concerns recognition and acceptance, as Augustine explains, rather than missing information in God. Gregory describes the present as the time for conversion and the closed door as the end of that opportunity. The appeal from outside therefore cannot repair what was left undone merely by repeating the correct address. Public belonging, admirable discipline and other people's testimony are each insufficient to supply the missing response.\footnote{Augustine, Sermon 93, §§15--16; Gregory, \textit{Homiliae in Evangelia}, 12.4--6.}

The four senses distinguish what the two interpretations explain. Literally, the narrative concerns preparation before a delayed arrival and the exclusion that follows its neglect. Allegorically, Christ is the Bridegroom who gathers and judges his Church; Gregory's Emmaus explanation concerns this same Lord received as a guest. Morally, Chrysostom directs attention to help withheld, while Augustine and Gregory examine the purpose of works already performed. Anagogically, the judgment discloses both the life lived and the charity within it, with lasting life in Christ as the promised end.

The Prayer over the Offerings concerns the Passion celebrated sacramentally and received inwardly. Read with the first interpretation, that reception calls for costly help after the pattern of Christ's self-giving. Read with the second, it includes conversion of the purpose with which the person participates and acts. The outward celebration remains a sacramental gift; the inward response concerns how its recipient lives.

Thanksgiving after Communion then includes asking the Spirit to preserve integrity. One act of generosity does not settle every future demand, and the correction of one vain purpose does not make renewed help unnecessary. The assembly asks for endurance as recipients who still depend on God. The two interpretations conclude in compatible responses: give the assistance that another person needs, and seek the charity by which that assistance is offered to God and neighbor.
